\section{Introduction}\label{sec:introduction}

\IEEEPARstart{T}{he} transition toward inverter-dominated power systems is
reducing grid inertia and accelerating frequency dynamics. Frequency
deviations and Rate of Change of Frequency (RoCoF) can exceed protection
thresholds within milliseconds~\cite{Pinheiro2025,WeakGrids2023}. Unlike
synchronous machines, Inverter-Based Resources (IBRs) regulate power
electronically, decoupling the primary energy source from electrical
frequency and increasing the impact of harmonics, phase discontinuities,
and nonstationary disturbances on protection and dynamic security
assessment~\cite{Bernal2024,Zoghby2024,Pinheiro2025}. In IBR-dominated
grids, phase discontinuities can induce cycle-slips in phase-locked loops,
sustained RoCoF can exceed the tracking capacity of multi-cycle window
methods, and converter-driven harmonic injection can cause spectral leakage
in DFT-based estimators. These failure modes may occur simultaneously, yet
existing benchmarks typically assess them in isolation~\cite{Ting2022},
leaving a gap for protection-class applications.

Each estimator family balances latency and robustness differently.
SRF-PLL is widely deployed but susceptible to cycle-slip under abrupt
phase discontinuities~\cite{EPRIPaper,Roscoe2019}. Window-based methods
(IpDFT, TFT) provide harmonic rejection at the cost of structural
observation latency~\cite{Belega2014,PlatasGarza2011}. Model-based filters
(EKF, UKF) improve transient tracking~\cite{Ferrero2020EKF,Chamorro2021},
while data-driven approaches offer potential accuracy gains at higher
computational cost~\cite{Ting2022,Zelaya2025}.

This paper presents a stress-oriented benchmark of twelve estimators from
five algorithmic families, evaluated in a local, non-time-synchronized
relay-class configuration under five scenarios: (A)--(C) follow the
IEC/IEEE~60255-118-1 test philosophy; (D)--(E) introduce IBR-specific
composite disturbances beyond the standard scope. The main contributions are:
(1)~a benchmark of twelve estimators from five algorithmic families under
five stress scenarios, including two composite IBR-specific scenarios
(D--E) not covered by IEC/IEEE~60255-118-1;
(2)~scenario-wise grid-search tuning with dual-rate simulation
(1\,MHz / 10\,kHz), establishing the performance ceiling per
estimator--scenario pair; and
(3)~a comparative analysis of RMSE, peak error, trip-risk, and
computational cost, providing selection guidance for relay-class
deployment in low-inertia grids.